\chapter{Amarás o Senhor, teu Deus, de todo o coração, de toda a alma e de todo o entendimento}

Jesus resumiu os deveres do homem para com Deus nestas palavras: «Amarás o Senhor teu Deus com todo o teu coração, com toda a tua alma, com toda a tua mente» (Mt 22,37), eco imediato do apelo solene do Deuteronómio: «Escuta, Israel: o Senhor nosso Deus é o único» (Dt 6,4). Deus foi o primeiro a amar; o amor do Deus único é lembrado na primeira das «dez palavras», e em seguida os mandamentos explicitam a resposta de amor que o homem é chamado a dar ao seu Deus. Os parágrafos 2083 a 2195 do Catecismo da Igreja Católica contemplam os três primeiros mandamentos do Decálogo --- os que dizem respeito mais diretamente ao amor de Deus: o Primeiro Mandamento, que chama o homem a crer, esperar e amar a Deus sobre todas as coisas; o Segundo, que exige o respeito pelo santo nome do Senhor; e o Terceiro, que manda santificar o Dia do Senhor. Juntos, estes três mandamentos regulam a relação pessoal e social do homem com Deus, da adoração interior às expressões externas do culto e do repouso sagrado.

\section{O Primeiro Mandamento (2083--2141)}

Deus dá-Se a conhecer lembrando a sua ação omnipotente, benevolente e libertadora: «Sou Eu que te tirei da terra do Egipto, dessa casa da escravidão» (Dt 5,6). O primeiro apelo e a justa exigência de Deus é que o homem O acolha e O adore: «Ao Senhor, teu Deus, adorarás, a Ele servirás; não ireis atrás de outras divindades» (Dt 6,13-14). O primeiro dos preceitos abrange a fé, a esperança e a caridade --- «pois quem diz Deus diz um ser constante, imutável, sempre o mesmo, fiel, perfeitamente justo», que merece plena confiança; «todo-poderoso, clemente, infinitamente propenso a bem-fazer», que merece toda a esperança; e gerador de amor, pois derramou sobre nós tesouros de bondade e ternura (\textit{Catecismo Romano}). A revelação da vocação e da verdade do homem está ligada à revelação de Deus: o Deus único e verdadeiro revela, antes de mais, a sua glória; o homem tem a vocação de manifestar Deus pelo seu agir, em conformidade com a sua criação «à imagem e semelhança de Deus» (Gn 1,26).

A nossa vida moral tem a sua fonte na fé em Deus, que nos revela o seu amor; o nosso dever para com Deus é crer n'Ele e dar testemunho d'Ele. O primeiro mandamento ordena que alimentemos e guardemos com prudência a nossa fé, rejeitando tudo quanto a ela se opõe. Pode-se pecar contra a fé de vários modos: a dúvida voluntária negligencia ou recusa ter por verdadeiro o que Deus revelou; a incredulidade é o desprezo da verdade revelada ou a recusa voluntária de lhe prestar assentimento; a heresia é «a negação pertinaz, depois de recebido o Batismo, de alguma verdade que se deve crer com fé divina e católica»; a apostasia é «o repúdio total da fé cristã»; e o cisma, «a recusa da sujeição ao Sumo Pontífice ou da comunhão com os membros da Igreja que lhe estão sujeitos» (CIC, can. 751). O primeiro mandamento inclui igualmente os pecados contra a esperança: o desespero, pelo qual o homem deixa de esperar de Deus a salvação pessoal ou o perdão dos seus pecados, opondo-se à bondade, à justiça e à misericórdia divinas; e a presunção, pela qual se confia excessivamente nas próprias capacidades, esperando salvar-se sem o socorro de Deus, ou se presume da omnipotência divina, esperando o perdão sem conversão. Finalmente, pode-se pecar contra o amor de Deus de diversas maneiras: a indiferença, que descuida ou recusa a consideração da caridade divina; a ingratidão, que não a reconhece nem a retribui amor com amor; a tibieza, hesitação em corresponder ao amor divino, que pode implicar a recusa de se entregar ao movimento da caridade; a acédia ou preguiça espiritual, que recusa a alegria que vem de Deus e aborrece o \textbf{bem} divino; e o ódio a Deus, que nasce do orgulho, nega a sua bondade e ousa amaldiçoá-Lo.

As virtudes teologais da fé, da esperança e da caridade informam e vivificam as virtudes morais; a virtude da religião dispõe-nos a prestar a Deus o que Lhe é devido com toda a justiça. Os atos próprios desta virtude --- a adoração, a oração, o sacrifício, as promessas e os votos --- exprimem a obediência ao primeiro mandamento. A adoração é o primeiro ato da virtude da religião: «adorar a Deus é reconhecê-Lo como tal, Criador e Salvador, Senhor e Dono de tudo quanto existe, amor infinito e misericordioso». O voto, «promessa deliberada e livre feita a Deus de um \textbf{bem} possível e melhor, deve cumprir-se por virtude da religião» (CIC, can. 1191 §1); é impossível alguém apropriar-se dos \textbf{bens} espirituais, que têm a sua fonte em Deus e só d'Ele se podem receber gratuitamente --- daí a proibição da simonia. O dever de prestar a Deus um culto autêntico diz respeito ao homem individual e socialmente: é dever social dos cristãos respeitar e despertar em cada homem o amor da verdade e do \textbf{bem}. O primeiro mandamento proíbe também a superstição e a idolatria --- que não diz respeito apenas aos falsos cultos do paganismo, mas é uma tentação constante: há idolatria desde o momento em que o homem honra e reverencia uma criatura em lugar de Deus (o poder, o prazer, o dinheiro, o Estado) ---; a adivinhação e a magia, que encerram uma vontade de dominar o tempo, a história e os homens, em contradição com o respeito que devemos a Deus e só a Ele; e os pecados de irreligião: tentar a Deus por palavras ou atos, o sacrilégio e a simonia. Na medida em que rejeita ou recusa a existência de Deus, o ateísmo é um pecado contra a virtude da religião; e o agnosticismo, com frequência, equivale a um ateísmo prático. O culto cristão das imagens sagradas, por sua vez, funda-se no mistério da encarnação do Verbo e não é contrário ao primeiro mandamento: «a honra prestada a uma imagem remonta ao modelo original» (II Concílio de Niceia, ano de 787).

\begin{description}

    \item[Pecados contra a virtude da fé:]
          \begin{itemize}
              \item \textit{Dúvida voluntária:} negligenciar ou recusar ter por verdadeiro o que Deus revelou e a Igreja propõe crer; quando deliberadamente cultivada, pode levar à cegueira do espírito.
              \item \textit{Incredulidade:} desprezo da verdade revelada ou recusa voluntária de lhe prestar assentimento.
              \item \textit{Heresia:} negação pertinaz, depois do Batismo, de alguma verdade que se deve crer com fé divina e católica, ou dúvida pertinaz acerca da mesma.
              \item \textit{Apostasia:} repúdio total da fé cristã.
              \item \textit{Cisma:} recusa da sujeição ao Sumo Pontífice ou da comunhão com os membros da Igreja a ele sujeitos.
          \end{itemize}

    \item[Pecados contra a virtude da esperança:]
          \begin{itemize}
              \item \textit{Desespero:} cessar de esperar de Deus a salvação pessoal, os socorros para a atingir ou o perdão dos pecados; opõe-se à bondade, à justiça e à misericórdia divinas.
              \item \textit{Presunção:} esperar salvar-se sem a ajuda do Alto (presunção das próprias capacidades), ou esperar o perdão sem conversão e a glória sem merecê-la (presunção da misericórdia divina).
          \end{itemize}

    \item[Pecados contra a virtude da caridade:]
          \begin{itemize}
              \item \textit{Indiferença:} descuidar ou recusar a consideração da caridade divina; desconhecer-lhe o cuidado preveniente e negar-lhe a força.
              \item \textit{Ingratidão:} não reconhecer, por desleixo ou recusa formal, a caridade divina; não retribuir amor com amor.
              \item \textit{Tibieza:} hesitação ou negligência em corresponder ao amor divino; pode implicar a recusa de se entregar ao movimento da caridade.
              \item \textit{Acédia (preguiça espiritual):} recusar a alegria que vem de Deus e aborrecer o \textbf{bem} divino.
              \item \textit{Ódio a Deus:} nasce do orgulho; opõe-se ao amor de Deus, nega a sua bondade e ousa amaldiçoá-Lo como Aquele que proíbe o pecado e lhe inflige o castigo.
          \end{itemize}

    \item[Outros pecados contra o Primeiro Mandamento:]
          \begin{itemize}
              \item \textit{Superstição:} atribuir uma eficácia de algum modo mágica a certas práticas ou aos sinais sacramentais, independentemente das disposições interiores que exigem.
              \item \textit{Idolatria:} honrar e reverenciar uma criatura (poder, prazer, dinheiro, Estado, etc.) em lugar de Deus; é incompatível com a comunhão divina.
              \item \textit{Adivinhação:} recorrer a horóscopos, astrologia, quiromância, vidência, médiuns ou outras práticas supostamente «reveladoras» do futuro; encerra uma vontade de dominar o tempo, a história e os homens.
              \item \textit{Magia e feitiçaria:} pretender domesticar poderes ocultos para os pôr ao seu serviço ou obter poder sobrenatural sobre o próximo; gravemente contrário à virtude de religião.
              \item \textit{Tentar a Deus:} pôr à prova, por palavras ou atos, a bondade e a omnipotência de Deus; implica sempre uma dúvida relativamente ao seu amor, à sua providência e ao seu poder.
              \item \textit{Sacrilégio:} profanar ou tratar indignamente os sacramentos, as ações litúrgicas, as pessoas, as coisas e os lugares consagrados a Deus; é pecado grave, sobretudo quando cometido contra a Eucaristia.
              \item \textit{Simonia:} comprar ou vender as realidades espirituais, que têm a sua fonte em Deus e só d'Ele se podem receber gratuitamente.
              \item \textit{Ateísmo:} rejeitar ou recusar a existência de Deus; é pecado contra a virtude da religião.
          \end{itemize}

\end{description}

\section{O Segundo Mandamento (2142--2166)}

O segundo mandamento manda respeitar o nome do Senhor. Depende, como o primeiro mandamento, da virtude da religião, e regula de modo particular o uso da palavra nas coisas santas. Entre todas as palavras da Revelação, há uma singular que é a revelação do nome de Deus: Deus confia o seu nome a quem crê n'Ele, revela-Se no seu mistério pessoal, e o dom do nome é da ordem da confidência e da intimidade. «O nome do Senhor é Santo»; por isso o homem não pode abusar dele: deve guardá-lo na memória «num silêncio de adoração amorosa» e não o empregará nas suas próprias palavras senão para o bendizer, louvar e glorificar. A deferência para com o nome de Deus exprime a deferência devida ao mistério do próprio Deus e a toda a realidade sagrada que esse nome evoca; o sentido do sagrado deriva da virtude da religião. O fiel deve dar testemunho do nome do Senhor, confessando a sua fé sem ceder ao medo; a pregação e a catequese devem estar compenetradas de adoração e respeito pelo nome de nosso Senhor Jesus Cristo.

O segundo mandamento proíbe o abuso do nome de Deus, isto é, todo o uso inconveniente do nome de Deus, de Jesus Cristo, da Virgem Maria e de todos os santos. A blasfémia opõe-se diretamente ao segundo mandamento: consiste em proferir contra Deus --- interior ou exteriormente --- palavras de ódio, de censura, de desafio; dizer mal de Deus; faltar-Lhe ao respeito; abusar do seu nome. A proibição estende-se às palavras contra a Igreja de Cristo, contra os santos e contra as coisas sagradas; é também blasfematório recorrer ao nome de Deus para justificar práticas criminosas, reduzir povos à escravidão, torturar ou condenar à morte. «A blasfémia é, em si mesma, pecado grave» (CIC, can. 1369). O segundo mandamento proíbe igualmente o juramento falso: fazer um juramento é tomar a Deus como testemunha do que se afirma; o juramento falso invoca Deus como testemunha de uma mentira. Comete perjúrio quem, sob juramento, faz uma promessa que não tem intenção de cumprir; o perjúrio constitui grave falta de respeito para com o Senhor de toda a palavra. Jesus expôs o segundo mandamento no Sermão da Montanha: «A vossa linguagem deve ser: "Sim, sim; Não, não". O que passa disto vem do Maligno» (Mt 5,37), ensinando que a presença de Deus e da sua verdade deve ser honrada em toda a palavra; a santidade do nome de Deus exige que não se recorra a ele por questões fúteis.

No Batismo, o cristão recebe o seu nome na Igreja. Pode ser o de um santo --- discípulo que levou uma vida de fidelidade exemplar ao seu Senhor --- cujo patrocínio oferece um modelo de caridade e assegura a sua intercessão; ou pode exprimir um mistério cristão ou uma virtude cristã: «procurem os pais, os padrinhos e o pároco que não se imponham nomes alheios ao sentir cristão» (CIC, can. 855). O cristão começa o seu dia, as suas orações e as suas atividades pelo sinal da cruz «em nome do Pai e do Filho e do Espírito Santo. Ámen»; o batizado consagra o dia à glória de Deus e este sinal fortalece-o nas tentações e nas dificuldades. O nome de todo o homem é sagrado: «Deus chama a cada um pelo seu nome» (cf.~Is 43,1; Jo 10,3); o nome é a imagem da pessoa e exige respeito, como sinal da dignidade de quem por ele se identifica. O nome recebido é, enfim, um nome de eternidade: «ao vencedor [...] dar-lhe-ei uma pedra na qual estará escrito um novo nome, que ninguém conhece, a não ser aquele que a recebe» (Ap 2,17).

\begin{description}

    \item[Pecados contra o Segundo Mandamento:]
          \begin{itemize}
              \item \textit{Blasfémia:} proferir contra Deus --- interior ou exteriormente --- palavras de ódio, censura ou desafio; dizer mal de Deus, faltar-Lhe ao respeito ou abusar do seu nome; estende-se às palavras contra a Igreja, os santos e as coisas sagradas; é pecado grave em si mesmo.
              \item \textit{Juras vãs:} invocar o nome de Deus sem intenção de blasfémia mas sem o respeito devido à sua grandeza e majestade; falta de reverência para com o Senhor.
              \item \textit{Juramento falso (perjúrio):} invocar Deus como testemunha de uma mentira, ou fazer uma promessa sob juramento sem intenção de a cumprir; constitui grave falta de respeito para com o Senhor de toda a palavra.
              \item \textit{Uso mágico do nome divino:} empregar o nome de Deus com finalidade mágica ou manipuladora, como se fosse um instrumento de poder ao serviço da vontade humana.
          \end{itemize}

\end{description}

\section{O Terceiro Mandamento (2168--2195)}

O terceiro mandamento do Decálogo refere-se à santificação do sábado: «O sétimo dia é um sábado: um descanso completo consagrado ao Senhor» (Ex 31,15). A Escritura fundamenta este mandamento tanto na memória da criação --- «porque em seis dias o Senhor fez o céu e a terra [...], mas ao sétimo dia descansou» (Ex 20,11) --- como no memorial da libertação de Israel da escravidão do Egipto: «Recorda-te de que foste escravo no país do Egipto [...]. É por isso que o Senhor, teu Deus, te ordenou que guardasses o dia de sábado» (Dt 5,15). Deus confiou a Israel o sábado como sinal da Aliança inviolável, santamente reservado ao louvor de Deus, da sua obra criadora e das suas ações salvíficas a favor do povo. O agir de Deus é o modelo do agir humano: se Deus «descansou» no sétimo dia, o homem deve também «descansar» e deixar que os outros, sobretudo os pobres, «tomem fôlego»; o sábado faz cessar os trabalhos quotidianos e é um dia de protesto contra as servidões do trabalho e o culto do dinheiro.

Jesus nunca violou a santidade do sábado, mas é com autoridade que lhe deu a interpretação autêntica: «O sábado foi feito para o homem e não o homem para o sábado» (Mc 2,27); cheio de compaixão, Cristo autorizou-Se a \textbf{fazer o bem} em dia de sábado, «a salvar uma vida antes que perdê-la» (cf.~Mc 3,4). Jesus ressuscitou de entre os mortos «no primeiro dia da semana» (Mc 16,2): enquanto «primeiro dia», o dia da ressurreição lembra a primeira criação; enquanto «oitavo dia», a seguir ao sábado, significa a nova criação inaugurada com a ressurreição de Cristo. O domingo tornou-se para os cristãos o dia do Senhor: «Reunimo-nos todos no dia do Sol, porque foi o primeiro dia em que Deus criou o mundo, mas também porque Jesus Cristo, nosso Salvador, nesse mesmo dia ressuscitou dos mortos» (São Justino, \textit{Apologia}). O domingo distingue-se expressamente do sábado, ao qual sucede cronologicamente em cada semana e cuja prescrição ritual substitui para os cristãos; realiza plenamente, na Páscoa de Cristo, a verdade espiritual do sábado judaico e anuncia o descanso eterno do homem em Deus. A celebração dominical da Eucaristia está no coração da vida da Igreja: «No domingo e nos outros dias festivos de preceito, os fiéis têm obrigação de participar na Missa» (CIC, can. 1247); os que deliberadamente faltam a esta obrigação cometem um pecado grave.

A instituição do Dia do Senhor contribui para que todos gozem de tempo suficiente de descanso e lazer, que lhes permita cultivar a vida familiar, cultural, social e religiosa. Aos domingos e festas de preceito, os fiéis abstenham-se dos trabalhos e negócios que impeçam o culto devido a Deus, a alegria própria do Dia do Senhor, a prática das obras de misericórdia ou o devido repouso do espírito e do corpo. O domingo é também um tempo de obras boas e de serviços humildes aos doentes, enfermos e idosos; um tempo de reflexão, silêncio, cultura e meditação, que favorecem o crescimento da vida interior e cristã. Santificar os domingos e festas exige um esforço comunitário: cada cristão deve evitar impor a outrem, sem necessidade, o que possa impedi-lo de guardar o Dia do Senhor; os poderes públicos têm a obrigação de assegurar aos cidadãos um tempo destinado ao repouso e ao culto divino; e os patrões têm obrigação análoga para com os seus empregados. No respeito pela liberdade religiosa e pelo \textbf{bem} comum de todos, os cristãos devem esforçar-se pelo reconhecimento dos domingos e dias santos da Igreja como dias feriados legais, dando a todos o exemplo público de oração, respeito e alegria, e defendendo as suas tradições como uma contribuição preciosa para a vida espiritual da sociedade humana.
